\documentclass[11pt]{article}

\usepackage[margin=1in]{geometry}
\usepackage{amsmath}
\usepackage{amssymb}
\usepackage{booktabs}
\usepackage{hyperref}
\usepackage{xcolor}

\hypersetup{
  colorlinks=true,
  linkcolor=blue,
  urlcolor=blue,
  citecolor=blue
}

\title{Derivation of the Ultrasound RBC Sensor SNR}
\author{Grand Unified Theory of Imaging (GUTI)}
\date{\today}

\begin{document}
\maketitle

\section*{Summary}

The ultrasound summary table reports an amplitude SNR of approximately
\[
  \mathrm{SNR}_{\mathrm{sensor}} = 0.1445.
\]
This number is not the SNR of a single isolated source--sensor pair.  It is the
RMS output amplitude of the full cone-scaled RBC pressure operator, restricted
to one 2 MHz axial range slice and multiplied by a 1\% CBV-variability factor,
divided by the assumed scalar sensor noise:
\[
  \mathrm{SNR}_{\mathrm{sensor}}
  =
  \frac{0.00072268~\mathrm{Pa}}{0.005~\mathrm{Pa}}
  =
  0.1445.
\]

\section*{Assumptions and References}

The numerical assumptions are taken from the repository implementation and
stored result metadata:
\begin{itemize}
  \item RBC backscatter defaults and metadata fields:
    \texttt{guti/modalities/us/analytical.py}.
  \item Modality-summary constants and SNR calculation:
    \texttt{scripts/make\_modality\_noise\_summary.py}.
  \item Stored 6000-sensor cone-scaled RBC result:
    \texttt{results/variants/us\_free\_field\_analytical\_50khz\_128k\_rbc\_cone\_slice\_cbv1pct\_20260602/json/001\_50khz\_128000src\_6000sensors.json}.
  \item Generated modality summary:
    \texttt{results/modality\_correlated\_noise\_summary/README.md}.
\end{itemize}

The key scalar assumptions are:
\[
\begin{aligned}
  f_{\mathrm{RBC}} &= 2~\mathrm{MHz},&
  f_{\mathrm{ref}} &= 50~\mathrm{kHz},&
  P_{\mathrm{external}} &= 1~\mathrm{MPa},\\
  \mathrm{BSC}_{10\mathrm{MHz}} &= 3\times 10^{-5}~\mathrm{cm}^{-1}\mathrm{sr}^{-1},&
  \mathrm{CBV} &= 0.03,&
  V_{\mathrm{voxel}} &= 24~\mathrm{mm}^{3},\\
  r_{\mathrm{ref}} &= 0.10~\mathrm{m},&
  T_{\mathrm{in}} &= 0.5,&
  T_{\mathrm{out}} &= 0.1.
\end{aligned}
\]

The \(24~\mathrm{mm}^{3}\) voxel volume is the backscatter/source-cell volume
used in the 50 kHz reference simulation basis.  It is not a claim that the
2 MHz image-resolution voxel has volume \(24~\mathrm{mm}^{3}\).  The later
2 MHz extrapolation uses the explicit spatial-mode scaling
\[
  \left(\frac{f_{\mathrm{RBC}}}{f_{\mathrm{ref}}}\right)^3
  =
  \left(\frac{2~\mathrm{MHz}}{50~\mathrm{kHz}}\right)^3
  =
  64{,}000.
\]

\section*{RBC Backscatter Model}

The volume-backscatter model converts a blood backscatter coefficient into a
pressure-amplitude transfer ratio.  The blood backscatter coefficient is
scaled from the 10 MHz reference value as
\[
  \mathrm{BSC}_{\mathrm{blood}}(f)
  =
  \mathrm{BSC}_{10\mathrm{MHz}}
  \left(\frac{f}{10~\mathrm{MHz}}\right)^4.
\]
At \(2~\mathrm{MHz}\),
\[
  \mathrm{BSC}_{\mathrm{blood}}(2~\mathrm{MHz})
  =
  3\times 10^{-5}
  \left(\frac{2}{10}\right)^4
  =
  4.8\times 10^{-8}~\mathrm{cm}^{-1}\mathrm{sr}^{-1}.
\]

Multiplying by cerebral blood volume gives the volume backscatter density
used for brain blood:
\[
  \eta
  =
  \mathrm{CBV}\cdot \mathrm{BSC}_{\mathrm{blood}}
  =
  0.03\cdot 4.8\times 10^{-8}
  =
  1.44\times 10^{-9}~\mathrm{cm}^{-1}\mathrm{sr}^{-1}.
\]
Converting from \(\mathrm{cm}^{-1}\) to \(\mathrm{m}^{-1}\):
\[
  \eta = 1.44\times 10^{-7}~\mathrm{m}^{-1}\mathrm{sr}^{-1}.
\]

The pairwise pressure model is
\[
  p_{ij}
  =
  P_{\mathrm{external}}\,
  T_i T_j
  \frac{\sqrt{\eta V_{\mathrm{voxel}}}}{r_{ij}}.
\]
The square root appears because the BSC gives an intensity-like scattering
fraction, while the channel model uses pressure amplitude.

\section*{Reference Pair Pressures}

For the reference range \(r=0.10~\mathrm{m}\) and
\(V_{\mathrm{voxel}}=24~\mathrm{mm}^3=24\times 10^{-9}~\mathrm{m}^3\),
\[
  \frac{\sqrt{\eta V_{\mathrm{voxel}}}}{r}
  =
  \frac{\sqrt{(1.44\times 10^{-7})(24\times 10^{-9})}}{0.10}
  =
  5.8788\times 10^{-7}.
\]
With \(P_{\mathrm{external}}=10^6~\mathrm{Pa}\), the no-skull reference pressure is
\[
  p_{\mathrm{no\ skull}}
  =
  10^6 \cdot 5.8788\times 10^{-7}
  =
  0.58788~\mathrm{Pa}.
\]

The cone run applies one skull-transmission factor on transmit and one on
receive.  Thus the three reference pair amplitudes are
\[
\begin{aligned}
  p_{\mathrm{out,out}}
    &= 0.58788 \cdot 0.1 \cdot 0.1
     = 0.0058788~\mathrm{Pa}
     = 5.8788~\mathrm{mPa},\\
  p_{\mathrm{in,out}}
    &= 0.58788 \cdot 0.5 \cdot 0.1
     = 0.029394~\mathrm{Pa}
     = 29.394~\mathrm{mPa},\\
  p_{\mathrm{in,in}}
    &= 0.58788 \cdot 0.5 \cdot 0.5
     = 0.14697~\mathrm{Pa}
     = 146.97~\mathrm{mPa}.
\end{aligned}
\]

These are baseline pair pressures before the 1\% CBV-variability factor.  The
variable component used in the table is \(0.01\) times the pressure response.

\section*{From Pair Pressures to the Table RMS Amplitude}

The table uses the full cone-scaled operator, not just one reference pair.  The
stored 6000-sensor SLQ metadata reports
\[
  \|G\|_F^2 = 372368.726,\qquad n_{\mathrm{outputs}}=366000.
\]
The full-volume, unit-variance output RMS is therefore
\[
  A_{\mathrm{all}}
  =
  \sqrt{\frac{\|G\|_F^2}{n_{\mathrm{outputs}}}}
  =
  \sqrt{\frac{372368.726}{366000}}
  =
  1.00866~\mathrm{Pa}.
\]

The summary row then restricts the source power to one axial range slice for a
two-cycle, 2 MHz pulse.  With speed of sound \(c=1540~\mathrm{m/s}\) and depth
\(D=0.15~\mathrm{m}\),
\[
  f_{\mathrm{slice}}
  =
  \frac{c(2/f_{\mathrm{RBC}})/2}{D}
  =
  \frac{1540(2/2{,}000{,}000)/2}{0.15}
  =
  0.00513333.
\]

The displayed table amplitude is
\[
  A_{\mathrm{table}}
  =
  A_{\mathrm{all}}\sqrt{f_{\mathrm{slice}}}\cdot 0.01
  =
  1.00866\sqrt{0.00513333}\cdot 0.01
  =
  0.00072268~\mathrm{Pa}.
\]
Equivalently,
\[
  A_{\mathrm{table}} = 0.72268~\mathrm{mPa}.
\]

\section*{Sensor Noise and Final SNR}

The summary row assumes scalar iid pressure noise
\[
  \sigma_{\mathrm{sensor}}=5~\mathrm{mPa}=0.005~\mathrm{Pa}.
\]
The reported sensor SNR is therefore
\[
  \mathrm{SNR}_{\mathrm{sensor}}
  =
  \frac{A_{\mathrm{table}}}{\sigma_{\mathrm{sensor}}}
  =
  \frac{0.00072268}{0.005}
  =
  0.1445359.
\]

\section*{Interpretation}

The number \(0.1445\) is a table-level RMS amplitude SNR under the current
RBC-cone assumptions.  It combines:
\begin{enumerate}
  \item BSC-derived RBC pressure scaling at 2 MHz,
  \item cone-dependent skull transmission,
  \item the 50 kHz reference simulation basis using \(V=24~\mathrm{mm}^{3}\),
  \item restriction to one two-cycle 2 MHz axial slice,
  \item a 1\% CBV-variability amplitude factor, and
  \item fixed \(5~\mathrm{mPa}\) scalar sensor noise.
\end{enumerate}

It should not be read as the SNR of every individual voxel-to-sensor pair.
For example, the 10 cm reference pair pressures range from
\(5.88~\mathrm{mPa}\) outside/outside to \(146.97~\mathrm{mPa}\) inside/inside
before the 1\% variability factor.

\section*{Repository References}

\begin{thebibliography}{9}
\bibitem{analytical}
GUTI ultrasound analytical implementation,
\texttt{guti/modalities/us/analytical.py}.

\bibitem{summary-script}
GUTI modality summary generator,
\texttt{scripts/make\_modality\_noise\_summary.py}.

\bibitem{result-json}
GUTI 6000-sensor RBC cone slice 1\% CBV result metadata,
\texttt{results/variants/us\_free\_field\_analytical\_50khz\_128k\_rbc\_cone\_slice\_cbv1pct\_20260602/json/001\_50khz\_128000src\_6000sensors.json}.

\bibitem{summary-readme}
GUTI generated modality noise summary,
\texttt{results/modality\_correlated\_noise\_summary/README.md}.
\end{thebibliography}

\end{document}
